\section*{Chapitre \ref{ch:mx:beyond-almgren-chriss} --- Au-delà d'Almgren--Chriss}

\begin{solution}{exo:mx:beyond-almgren-chriss:1}
$\phi=e^{-\ln2\,(1/78)/0.2}=0.9565$.
\end{solution}

\begin{solution}{exo:mx:beyond-almgren-chriss:2}
Le coût comporte le terme $-\tau\alpha x$~: détenir des actions pendant que la dérive est positive rapporte la dérive, si bien qu'un vendeur conserve davantage de sa position lorsque $\alpha>0$
et vend moins~; lorsque $\alpha<0$, il vend plus vite.
\end{solution}

\begin{solution}{exo:mx:beyond-almgren-chriss:3}
$1/(\rho T+2)=1/12=0.0833$ de l'ordre dans chaque bloc (0.0876 sur la grille de 101 transactions).
\end{solution}

\begin{solution}{exo:mx:beyond-almgren-chriss:4}
Agir sur une prévision revient à trader plus vite ou plus lentement que le calendrier, ce qui coûte un impact temporaire~; or une prévision qui s'éteint en quelques pas modifie
peu le cours avant de disparaître, de sorte qu'il y a peu à capter pour l'impact dépensé~: 0.19 point de base pour un vingtième de journée et 0.5~USD par jour, 3.0 même
à 2~USD.
\end{solution}

\begin{solution}{exo:mx:beyond-almgren-chriss:5}
La stratégie change tellement de vitesse après les mouvements qu'elle paie un impact élevé par à-coups~: son coût moyen grimpe à 81.7 points de base, ce qui dépasse la valeur du risque
qu'elle supprime à cette aversion au risque (équivalent certain 103.6 contre 101.9).
\end{solution}

\begin{solution}{exo:mx:beyond-almgren-chriss:6}
Elle trade davantage lorsque le marché est liquide et attend lorsqu'il est sec (60\,000 contre 20\,000 actions à mi-ordre), payant moins souvent l'impact quadruple du régime sec~;
Almgren--Chriss trade selon le calendrier quel que soit le régime.
\end{solution}

\begin{solution}{exo:mx:beyond-almgren-chriss:7}
Agir sur l'observation bruitée comme si elle était exacte économise 8.3 points de base (erreur type 2.0) au lieu de 41~; en la pondérant par un demi, soit son rapport variance du signal sur variance totale
(l'estimateur bayésien de $\alpha$), on économise 20.3. Le bruit fait trader la politique sur des erreurs~; réduire la prévision en restitue la moitié de la valeur.
\end{solution}

\begin{solution}{exo:mx:beyond-almgren-chriss:8}
L'économie dépend de l'intensité et de la demi-vie du signal relativement à la durée et à la taille de chaque ordre, elle se mesure sur les ordres du backtest, et elle
suppose que le signal reste aussi bon hors échantillon et qu'il est observé sans bruit (exercice~7~: le bruit la ramène de 41 à 8 points de base)~; et le calendrier prend davantage de risque.
Il faut l'évaluer par type d'ordre, hors échantillon, en tenant compte du risque.
\end{solution}

\begin{solution}{pb:mx:beyond-almgren-chriss:1}
\textbf{1.} Minimiser $\E\sum(\eta u^2/\tau+\lambda\sigma^2\tau x^2-\tau\alpha x)$~; état $(x,\alpha)$.
\textbf{2.} Coût quadratique et dynamique linéaire~: la fonction valeur est quadratique, la politique linéaire~; récurrence de Riccati à rebours.
\textbf{3.} Avec $\phi=0$ et sans dérive, les gains reproduisent la trajectoire discrète d'Almgren--Chriss.
\textbf{4.} Demi-vie 0.05~: 0.19, 0.75, 2.98~; 0.2~: 1.1, 4.5, 17.2~; 1~: 3.1, 11.5, 41.0 points de base pour 0.5, 1 et 2~USD par jour.
\textbf{5.} Plus rapide lorsque le cours a bougé dans le sens favorable à l'ordre.
\textbf{6.} 60.9, 90.5, 101.9 points de base.
\textbf{7.} Moyennes 61.2, 62.2, 66.2, 81.7~; écarts types 88.7, 86.4, 80.4, 66.2~; équivalents certains 100.5, 99.4, 98.6, 103.6.
\textbf{8.} Après un gain, réduire le risque restant vaut un certain impact~; le calendrier statique ne peut pas le faire.
\textbf{9.} 0.0893 et 0.0834~; blocs de 0.0876 (formule 0.0833).
\textbf{10.} Une résilience de 10 par jour fait décroître l'impact rapidement par rapport à la journée, si bien qu'un étalement régulier est presque optimal.
\textbf{11.} Normal et sec (quatre fois l'impact), rester à 0.9 et 0.7~; 60\,000 contre 20\,000 actions à mi-ordre.
\textbf{12.} 73.2 contre 91.5 points de base~: 18.3 économisés.
\textbf{13.} Signal (jusqu'à 41 pour un signal fort et persistant), liquidité (18), impact transitoire (6.6~\% de l'impact), adaptation (3.2~\% de l'équivalent certain).
\textbf{14.} Le signal et les régimes de liquidité.
\textbf{15.} Elle l'augmente~: 212 contre 173 points de base d'écart type pour le signal le plus fort.
\textbf{16.} Scinder les ordres réels selon la valeur du signal à l'arrivée et comparer les calendriers avec et sans lui sur des ordres appariés, hors échantillon.
\textbf{17.} \emph{Résultat nommé}~: sur le bloc, une prévision de dérive de demi-vie d'un vingtième de journée économise de 0.19 à 3.0 points de base pour des écarts types de
0.5 à 2~USD par jour, et une prévision de demi-vie d'une journée économise de 3.1 à 41~; l'économie croît avec l'intensité et plus vite que linéairement avec la persistance.
\textbf{18.} Environ un cinquième de journée (1 à 17 points de base) pour ces intensités~; les signaux plus brefs valent peu.
\textbf{19.} Régimes de liquidité~: une économie de 18 points de base à partir de données que le desk voit en temps réel.
\textbf{20.} La chance de trader lorsque trader est bon marché et que le cours va dans le même sens, au prix d'un risque accru.
\end{solution}

\begin{solution}{iq:mx:beyond-almgren-chriss:1}
Vendre plus vite maintenant, avant la baisse anticipée, au prix d'un impact plus élevé~; de combien plus vite dépend de l'ampleur et de l'horizon de la prévision face à mon impact.
\iqlookfor{Direction~; l'arbitrage avec l'impact~; horizon.}
\end{solution}

\begin{solution}{iq:mx:beyond-almgren-chriss:2}
État (positions, signal), contrôle de la transaction~; coût égal à l'impact plus le risque moins la dérive multipliée par les positions~; dynamique linéaire avec la persistance d'Ornstein--Uhlenbeck~; résolution par Riccati.
\iqlookfor{État, coût, dynamique, méthode de résolution.}
\end{solution}

\begin{solution}{iq:mx:beyond-almgren-chriss:3}
Un premier bloc utilise le carnet tel qu'il est, la partie continue trade au rythme de reconstitution, et le dernier bloc utilise une liquidité dont aucune transaction ultérieure ne profitera.
\iqlookfor{Résilience~; effets de bord.}
\end{solution}

\begin{solution}{iq:mx:beyond-almgren-chriss:4}
Accélérer, au sens moyenne-variance~: le gain paie la réduction du risque (agressive dans la monnaie)~; modérément, car une forte adaptation coûte trop d'impact.
\iqlookfor{Direction~; modération.}
\end{solution}

\begin{solution}{iq:mx:beyond-almgren-chriss:5}
Mêmes ordres, mêmes trajectoires de marché (nombres aléatoires communs), nombreuses répétitions, coûts mesurés par rapport au même prix d'arrivée, on rapporte les moyennes avec leurs erreurs types et le
risque.
\iqlookfor{Nombres aléatoires communs~; erreurs~; risque.}
\end{solution}

\begin{solution}{iq:mx:beyond-almgren-chriss:6}
Estimer le pouvoir prédictif du signal sur les rendements aux horizons d'exécution, simuler la politique sur des ordres historiques avec le signal tel qu'il est connu à l'arrivée,
et comparer avec le calendrier sans lui, hors échantillon et net de l'impact supplémentaire.
\iqlookfor{Hors échantillon~; adéquation des horizons~; net de l'impact.}
\end{solution}
